\documentclass[11pt]{letter}
\usepackage[utf8]{inputenc}
\usepackage[top=2.5cm, bottom=2.5cm, right=3cm, left=3cm]{geometry}
\usepackage{url}

\signature{C. Ben Gibson\\
RAND Corporation\\
\texttt{cbengibson@gmail.com}}

\address{C. Ben Gibson\\
RAND Corporation\\
1776 Main Street\\
Santa Monica, CA 90401}

\begin{document}

\begin{letter}{The Editors\\
\textit{Social Networks}}

\opening{Dear Editors,}

I am pleased to submit for your consideration the manuscript ``These Four First:
A Robust Endogenous Baseline for Relational Event Sequences,'' which I am
submitting as an original research article to \textit{Social Networks}.

Applied researchers using relational event models (REMs) face a recurring
specification problem: a growing catalogue of endogenous statistics is available,
but no consensus exists on which subset any well-specified model must include.
Omitting a consistently important endogenous term is not a minor oversight ---
its signal redistributes across the remaining predictors, potentially biasing
coefficients on the exogenous covariates of substantive interest. Yet including
all available statistics risks convergence failures, overfitting, and tangled
interpretation. The field has lacked empirical guidance on where to start.

This paper provides that guidance. I apply greedy forward AIC selection to 217
directed interaction sequences spanning ten social domains --- private messaging,
email, phone communication, face-to-face interaction, gang violence, and
international conflict --- and identify which endogenous statistics are
\emph{consistently first-selected} regardless of domain. The answer is four
terms, and the gap to the fifth is sharp. IOviolation (participation-shift
violations), SndInitiation (prior sender initiation), Repetition (normalized
dyadic send-share), and RecRecencyRank (rank-based reverse-dyad recency) each
appear in the greedy top-four of over 55\% of domains (domain-weighted), with a
22-percentage-point drop to the next candidate. Together they recover
approximately 97\% of the AIC improvement achievable across the twelve-term
candidate set. Two path-independent sensitivity criteria --- univariate AIC rank
and take-one-out from the full joint model --- corroborate the greedy result.

Several findings are of broader theoretical interest. First, the two most
prominently discussed mechanisms in the network literature --- imitation
(preferential attachment) and triadic closure --- do not belong in this baseline:
their signal is almost entirely absorbed by forward-dyad history once
SndInitiation and Repetition enter the model. Second, the ecological split
between high- and low-visibility interaction settings is consequential:
imitation is a serious candidate for inclusion in high-visibility contexts
(40\% top-four frequency) but negligible in private communication (5\%). Third,
the consistent early selection of SndInitiation --- a binary absorbing state
partitioning the risk set into first-contact and return-contact events ---
motivates a separable modeling approach, pointing toward a productive direction
for future methodological work.

I believe this manuscript is well suited to \textit{Social Networks} for several
reasons. It speaks directly to the practice of network analysis across
substantive domains represented in the journal's readership. It engages
methodological literature the journal has published (including work on REM
specification, goodness of fit, and sampling), and its cross-domain empirical
scope is consistent with the journal's emphasis on general structural and
dynamic network phenomena. The finding that a compact four-term baseline spans
settings from classroom interaction to gang violence is, I hope, of genuine
interest to the journal's broad readership.

This manuscript has not been previously published and is not under review
elsewhere. The datasets used are drawn from previously published corpora; full
descriptions and citations are provided in the appendix.

Thank you for your consideration.

\closing{Sincerely,}

\end{letter}
\end{document}
